\subsection*{الف}
فرض می‌کنیم قبل از آزمایش، هر توپ با احتمال یکسان می‌تواند توپ غیرعادی باشد و همچنین غیرعادی بودن آن (سنگین‌تر/سبک‌تر) با احتمال برابر رخ می‌دهد؛ پس $24$ حالت هم‌احتمال داریم.

در وزن‌کشی اول، $4$ توپ بیرون می‌مانند، $4$ توپ در کفه‌ی چپ و $4$ توپ در کفه‌ی راست قرار می‌گیرند. خروجی‌های ممکن:
\[
O \in \{\text{تراز},\ \text{کفه‌ی چپ سنگین‌تر},\ \text{کفه‌ی راست سنگین‌تر}\}.
\]
احتمال هر خروجی:
\begin{flalign*}
P(\text{تراز})
&= P(\text{توپ غیرعادی در 4 توپ بیرون باشد})
= \frac{4}{12}=\frac{1}{3},\\
P(\text{چپ سنگین‌تر})
&= P(\text{توپ غیرعادی سنگین و در چپ}) + P(\text{توپ غیرعادی سبک و در راست}) \\
&= 4\left(\frac{1}{12}\cdot\frac{1}{2}\right)+4\left(\frac{1}{12}\cdot\frac{1}{2}\right)
= \frac{8}{24}=\frac{1}{3},\\
P(\text{راست سنگین‌تر})
&= \frac{1}{3}. &&
\end{flalign*}

بنابراین آنتروپی (با لگاریتم مبنای 2، بر حسب \lr{bit}) برابر است با:
\begin{flalign*}
H_1
&= -\sum_{o} P(o)\log_2 P(o)
= -3\left(\frac{1}{3}\log_2\frac{1}{3}\right)
= \log_2 3
\approx 1.585~\text{bit}. &&
\end{flalign*}

\subsection*{ب}
در این حالت همه‌ی توپ‌ها وزن می‌شوند؛ پس خروجی \lr{تراز} غیرممکن است و تنها دو خروجی داریم:
\[
O \in \{\text{چپ سنگین‌تر},\ \text{راست سنگین‌تر}\}.
\]
به دلیل تقارن:
\begin{flalign*}
P(\text{چپ سنگین‌تر})
&= 6\left(\frac{1}{12}\cdot\frac{1}{2}\right)+6\left(\frac{1}{12}\cdot\frac{1}{2}\right)
= \frac{12}{24}=\frac{1}{2},\\
P(\text{راست سنگین‌تر})&=\frac{1}{2}. &&
\end{flalign*}
پس:
\begin{flalign*}
H_2
&= -2\left(\frac{1}{2}\log_2\frac{1}{2}\right)
= 1~\text{bit}. &&
\end{flalign*}

\subsection*{ج}
\[
H_1=\log_2 3 \approx 1.585~\text{bit} \;>\; H_2=1~\text{bit}.
\]
در وزن‌کشی $4$-در-برابر-$4$ سه خروجی هم‌احتمال داریم (به‌خصوص حالت تراز که نشان می‌دهد توپ غیرعادی در 4 توپ کنارگذاشته است)، بنابراین عدم‌قطعیت خروجی بیشتر و در نتیجه اطلاعات مورد انتظار بیشتر از حالت $6$-در-برابر-$6$ است که فقط دو خروجی هم‌احتمال دارد.